\subsection{Centering the exponent on its full support}
\label{subsec:selected-exponent}

The assigned primes fix the column dimensions. In the arithmetic comparison each lift contributes exponents at every prime in its support. Centering its contribution requires accounting for all of them.

\begin{paramdefinition}[Word, lift and local exponents]\label{hyp:selected-exponent-data}
Fix a word \(u\) and a lift \(1\le j\le m\) at the point of
Section~\ref{subsec:selected-word-level}. Let \(T\) index the designated top weights and \(a_t=a_t(u)\) be the weight of the sequence at \(t\). Write \(S=\sum_{t\in T}a_t\),
\(L=\log b\), \(B_0=fmbL\), and \(N=N_m\). For the coefficient field
\(K\) put \(\omega_w=\log\mathrm Nw/[K:\mathbb Q]\) and
\(e_w=v_w(p)\) at \(w\mid p\), with \(v_w\) the integer valuation. Use the local exponent
decomposition of Section~\ref{sec:finite-valuations} and the
level at \(T_u\) of Proposition~\ref{prop:assigned-prime}.
\end{paramdefinition}
\begin{proposition}[Uniform centering]\label{prop:selected-exponent-centering}
For the word and lift of Definition~\ref{hyp:selected-exponent-data},
let \(\mathcal S\) contain both full fibres above \(g_0,g_1\);
allow partial support at the remaining primes. Then
\begin{align}
 \left|\sum_{w\in\mathcal S}2\omega_w e_w E_{u,j}(w)
               -S\log(g_0g_1)\right|&\ll B_0,\label{eq:selected-exponent-all}\\
 \sum_{\substack{w\in\mathcal S\\w\nmid g_0g_1}}
             2\omega_w e_w |E_{u,j}(w)|&\ll B_0.\label{eq:selected-exponent-away}
\end{align}
The constants and eventual threshold can be taken before the
point, word and lift are chosen. The denominator determining the finite support may vary. The local exponents \(E_{u,j}(w)\) are those of Section~\ref{sec:finite-valuations}. \label{cor:selected-exponent-1}
\end{proposition}

\begin{proof}
\proofstep{The two complete type fibres.}
Let \(g\) be either type prime and \(\pi\bar\pi=g\) its norm
factorization. The level at the prime assigned to the sequence is prime to \(g\) and the support
contains the entire fibre. Let \(H_\rho\) be the sum of the
partition rows of the jet at \(\rho\), and \(J\) the total Fock
degree of the lift. At \(w\mid g\) the entry is
\[
 v_g(D_{\Gamma,u})-J/2+
 V_w+c_g f(N^2+m^2\eta),\qquad
 V_w=\sum_\rho H_\rho\,v_w(\rho(\bar\pi))/e_w .
\]
Set \(\mu_w=e_w\omega_w\). Conjugation permutes the places over
\(g\), preserving \(\mu_w\). Since each CM embedding commutes with
conjugation and \(\pi\bar\pi=g\), the following identities hold even
at a place fixed by this permutation:
\[
 V_w+V_{\bar w}=\sum_\rho H_\rho,\qquad
 \sum_{w\mid g}\mu_w=\log g .
\]
Summing on the complete fibre gives
\[
 2\sum_{w\mid g}\mu_w V_w
   =\sum_{w\mid g}\mu_w(V_w+V_{\bar w})
   =(\log g)\sum_\rho H_\rho .
\]
The other terms of the entry are constant on the fibre. Thus
the fibre sum minus \(S\log g\) is exactly
\[
 \left(\sum_\rho H_\rho-J-S+2v_g(D_{\Gamma,u})
                  +2c_g f(N^2+m^2\eta)\right)\log g.
\]

\proofstep{The jet weights and Fock degree.}

Put \(e=w_{j\rho}\), and let \(A_\rho=a_t\) on the top indexed
by \(t\), zero elsewhere. The first \(k\) rows of the jet have the
form \(u_r+e\) with shifts between \(-N\) and zero. The \(k-2\)
interior head parameters lie in the prescribed interval of
size \(S_c\). The
first parameter equals \(A_\rho\) at a top and has size at most
\(3S_c\) elsewhere. The last row of the head is positive and
at most the second row. Each of the remaining
\((j-1)(t+1)\) entries is a difference of two offsets in the
same width window with its rank shift, of size at most \(2\eta\).
Here \(t\) in the ranks denotes the block-length parameter. It
follows that
\[
 |H_\rho-A_\rho-(k-2)e|
 \le 2e+3k(S_c+N)+2m(t+1)\eta .
\]

For the local Fock degree \(J_\rho\) sum the absolute values of the
two letter weights. The first \(k-1\) entries of the large block
are \(e\) plus a head weight and a half-sum shift bounded by
\(2N\). The other \(j\) entries and the \(m-j\) entries of the
small block involve \(e-B_h\) or its negative with a shift of size
at most \(2N\), where all \(B_h\) at this place lie within
\(2\eta\) of \(e\). Using \(e\ge0\) and
\(\big||q|-e\big|\le |q-e|\) gives
\[
 |J_\rho-(k-1)e|
 \le k(S_c+2N)+2m(\eta+2N).
\]
The degree formula uses both blocks with the small block in its
dual order, hence \(J=\sum_\rho J_\rho\).

There is exactly one
\(\rho\) for each designated top. Summing the last two inequalities
and using \(\sum A_\rho=S\) gives
\[
 \left|\sum_\rho H_\rho-J-S\right|
 \ll \sum_\rho w_{j\rho}+f\big(k(S_c+N)+m(t+1)\eta\big)\ll fb .
\]
Indeed \(e\le 3b\), \(kS_c,kN,m(t+1)\eta\le b\) on the chosen
size ranges. Also \(N^2+m^2\eta\ll mb\), \(\log g\ll L\), and
\(v_g(D_{\Gamma,u})\log g\le\log D_{\Gamma,u}\ll
fN\log(2fN)\) by the common Minkowski stabilizer bound. The last
ranks grow polynomially with \(x\) while \(b=\exp(x^\epsilon)\).
This proves a two-sided \(O(B_0)\) error on each type fibre.

\proofstep{The assigned prime and the determinant factor.}
At the assigned prime \(T_u\), the two orientations have equal
total local degree by conjugation. In the notation of the assigned-prime
term \eqref{eq:tag-w5} of Section~\ref{subsec:word-period}, the two entries are
\[
 E_{u,j}(w)=
 \begin{cases}
 2\sum_\rho W_\rho+fN+|T|(t+1)(m+1)/4,&\text{chosen orientation},\\
 0,&\text{conjugate orientation}.
 \end{cases}
\] In \(W_\rho\) take the maximum of the
absolute entries of the successive weights after subtracting the
integer shift \(\lfloor i(t+1)/2\rfloor\) at step \(i\), omitting entry one at a top.
In the original head the remaining parameters have size
\(O(S_c+b')\), including the negative last parameter at a top.

In the appended blocks the parameters are \(-w_{h\rho}\), of
size at most \(3b\). Passing from these parameters to the
weight and removing the common shift costs \(O(N)\) in each
coordinate. Thus \(W_\rho\ll b\). Since \(\log T_u\ll mL\),
the fibre costs \(O(B_0)\). For the determinant term on the other fibres it suffices to take
absolute values and extend the sum to every rational prime. The bound is
\[
            2\sum_p v_p(D_{\Gamma,u})\log p=2\log D_{\Gamma,u}.
\]

\proofstep{The discriminant prime.}

In Proposition~\ref{prop:unserved-discriminant-entries}, let
\(A_t^\pm\) be the two column counts, \((A_t^-)'\) the length of
the second minus row, and \(H_{(-,t)}\) the total minus degree.
The discriminant calculation gives
\[
 c_t(u)=
 \begin{cases}
 H_{(-,t)},& t\text{ off the tops},\\
 H_{(-,t)}-A_t^-+(A_t^-)' + A_t^+,&t\text{ at a top}.
 \end{cases}
\]
Thus at a top only the excess of the first minus row over the
second can carry the exceptional parameter. At the minus embeddings the weights are in
the smaller range. The plus embeddings lie outside the top index set.
The minus chain weight is at most \(9b'/4\) and the plus
weight at most \(b+\eta+1\). Off the tops the entire minus degree
is \(O(kb')\) by the head and tail bounds above, using
\[
 kS_c+kN+m(t+1)\eta\ll kb'.
\]

At a top first subtract the
difference of the first two minus rows. This cancels the exceptional
\(a_t\), since the weight in the second row lies in the interval
for the interior head weights. Restoring the first
plus row costs at most \(b+O(S_c+\eta+N)\). Positivity and
monotonicity of the rows show \(c_t\ge0\) in both cases. There are \(f/2\) pairs in \(F_d\), and the smaller weight in each
pair is at most \(9b'/4\). Thus
\[
 0\le -F_d\le \frac{k-1}{2}\,\frac f2\,\frac{9b'}4 .
\]
The non-top terms in \(\sum_t c_t(u)\) cost \(O(fkb')\);
the top terms add \(O(|T|b)\) by the cancellation just obtained. Thus
\[
 |F_d+\tfrac12\max_u\sum_t c_t(u)|
          \ll |T|b+fkb'\ll fb .
\]
Here \(kb'\le b\). Multiplication by \(2\log d\le2L\) suffices.

\proofstep{Central and remaining inert fibres.}

The marked-prime contribution does not exhaust the exponent. The central and remaining inert fibres must be included before the total can be centered on its support.


The inert terms vanish above \(X''\). At a central residue prime \(p\le X''\), the entry is
\[
 3fN^2\log_p a+\frac{\sum_\rho w_{j\rho}}{p^2-1}.
\]
Weight by \(\log p\) and sum. For the set \(\mathcal P_c\) of central primes, \(\log_p a\log p=\log a\)
and \(\#\mathcal P_c\le X''\). Hence the sum is at most
\[
 3fN^2 X''\log a+
  \Big(\sum_\rho w_{j\rho}\Big)\sum_{n\ge2}\frac{\log n}{n^2-1}.
\]
The series converges by comparison with \((\log n)/n^2\). Since \(X''N^2\le b\) and
\(\log a\ll_\epsilon L\), the total is \(O_\epsilon(fbL)\).

At any other inert prime separate the unbarred orbits without a
signal (class A) from the remaining indices (classes B and C).
If \(J_i\) denotes the orbit objective starting at \(i\), the
quantity that must be bounded for class A is
\[
 A_p=\left(\frac1{n_p}\sum_{i\in {\rm A}}J_i\right)_+ .
\]
Taking the positive part after the sum preserves the common period
normalization of the local entry. The B and C contributions are nonnegative
for the chosen local constants. Their sum with the positive part gives the entry after removing
the exceptional contribution. Its cubic weight and digit cost
are estimated separately below.

To see this, test the
orbit minimum with zero. Its upper bound at a starting index is
\(p/(p^2-1)\) times the sum of \((k-1)w\) over its \(n_p\)
steps. Without the signal an unbarred starting index stays
unbarred, so \(w\le 3(b'+\eta)\) there, and
\(k(b'+\eta)\le mb\). Since there are at most \(f\) starts, the positive part of
the mean of the A objectives is at most
\[
 \frac{f}{n_p}\,\frac{p}{p^2-1}\,n_p\,3(k-1)(b'+\eta)
       \ll fmb/p .
\]

For B and C, the top count costs at most
\(fp(b+\eta+1+3S_c)/(p^2-1)\) and the sum of weight
terms at most \(12\sum w/p\ll fb/p\).
The digit and rank prices cost \(O(f\eta L)\) per prime
after multiplication by \(\log p\): the rank products
are at most \(N^2\le\eta\), and \(\log(2a),\log(4n_pNa)=O(L)\).
The squared base-\(p\) logarithm costs \(O(L^2/\log p)\),
also \(O(\eta L)\) there.
By Chebyshev, each dyadic interval below \(X''\) contributes
\(O(1)\) to \(\sum_{p\le X''}\log p/p\). Thus the weight
and digit prices total
\[
 O\left(fmb\sum_{p\le X''}\frac{\log p}{p}
                       +f\eta L X''\right)=O(fmbL),
\]
using \(X''\eta\le b\).

There is one further cubic weight sum. It occurs only at a
residue prime with nonzero Frobenius component of odd order, and
contributes \(O(fkb\log p/p^3)\) and a digit price already covered
above. The projection of this component to the dyadic group
\((\mathbb Z/2^{e_2+2}\mathbb Z)^\times/\{\pm1\}\) is trivial.
Thus, with \(H=2^{e_2+2}\), any such prime satisfies
\(p\equiv\pm1\pmod H\). Writing \(p=hH\pm1\) gives \(h\ge1\) for
either sign. Since \(hH-1\ge hH/2\),
\[
 \sum_{p\ {\rm in\ this\ class}}\frac{\log p}{p^3}
 \ll \frac1{H^3}\sum_{h\ge1}\frac{\log(2H)+\log h}{h^3}
 \ll\frac{\log(2H)}{H^3}.
\]
The bounds \(H\ge x^2\), \(k\le x^5\) imply
\(k\log(2H)/H^3=O(1)\), completing the inert prime sum.

These terms exhaust the off-type definition. On a subset of the
fibre of \(p\) the normalized local degrees sum to at most \(\log p\).
This proves \eqref{eq:selected-exponent-away}, even when the
support there is a proper subset. Adding the two complete type
fibres gives \eqref{eq:selected-exponent-all}. Summing over the
\(m\) lifts gives \(|E_u-mS\log(g_0g_1)|\ll mB_0\) on precisely
the words counted in Section~\ref{subsec:selected-word-level}.
\label{cor:selected-exponent-2}
\end{proof}
